\section{Future Directions}
\label{sec:future}

\Cref{sec:discuss} read the architecture documented in \Cref{sec:arch,sec:turn,sec:auth,sec:ext,sec:context,sec:subagent,sec:persist} as a coherent design point and surfaced the tensions, trade-offs, and near-horizon directions that design point implies. This section steps beyond the architecture itself to record six open questions that \Cref{sec:discuss:future} partially names and that a growing external literature has sharpened enough to state concretely. These questions cover the paper's five-value framework (\Cref{sec:values:five}) and the long-term developer-capability question introduced in \Cref{sec:values:lens}. They concern governance constraints on the Authority hierarchy (\Cref{sec:future:governance}), the observability-evaluation gap on the Safety side (\Cref{sec:future:evaluation}), cross-session persistence on the Reliability side (\Cref{sec:future:persistence}), four extensions of the Capability frontier (\Cref{sec:future:boundary}), horizon scaling as a distinct axis of Reliable Execution (\Cref{sec:future:horizon}), and long-term developer capability reframed as a design question (\Cref{sec:future:lens}). Consistent with \Cref{sec:discuss:future}'s framing, each question is posed in the form \emph{whether}/\emph{how}/\emph{which}. Specific mechanism choices are named when the cited sources name them and left open otherwise.

\subsection{Silent Failure and the Observability-Evaluation Gap}
\label{sec:future:evaluation}

The observability-evaluation gap reported in \Cref{sec:discuss:future} could reflect a missing tooling layer, a missing evaluation interface inside the harness, or a model-capability ceiling. The sources there do not adjudicate. How the silent-mistake failure mode noted in that paragraph should be surfaced is therefore an architectural question for the harness rather than a capability question for the model. Recent empirical work characterises the gap at several resolutions. \citet{cemri2025massfail} catalogue fourteen failure modes spanning system-design issues, inter-agent misalignment, and task verification. \citet{pathak2025silentmultiagent} build a benchmark of agent trajectories specifically for anomaly detection in traces. \citet{yao2024taubench} expose consistency gaps via the $\text{pass}^k$ metric, the probability that all $k$ independent trials succeed. \citet{kapoor2024agentsthatmatter} argue that current agent benchmarks lack holdouts and cost controls, limiting what observability can actually diagnose.

Against the permission pipeline and tool-orchestration layers analysed in \Cref{sec:auth,sec:turn}, two architectural questions remain open. First, whether the scaffolding the paper cites from \citet{anthropic2026harness} (generator-evaluator separation, sprint contracts, post-hoc checks, building on \citet{madaan2023selfrefine}'s self-refine pattern) belongs inside the harness (e.g., as additional hook events alongside the 27 documented in \Cref{sec:ext}) or outside it as a separate evaluation layer is not settled by the cited sources. Second, whether the existing hook pipeline of \Cref{sec:ext} can host such scaffolding within its current context-cost envelope is a further open question. The observation that closing this gap ``likely requires additional scaffolding \ldots rather than model improvements alone'' (\Cref{sec:discuss:future}) locates the open work at the harness layer.

\subsection{Persistence: Memory and Longitudinal Colleague Relationships}
\label{sec:future:persistence}

Whether agent state and the human-agent working relationship should persist across sessions, and in what form, is treated by the paper at two distinct layers today. \Cref{sec:context} documents the four-level CLAUDE.md hierarchy and auto memory. \Cref{sec:persist} documents mostly-append-only JSONL transcripts whose session-scoped permissions resume does not restore, apart from explicit cleanup rewrites. What belongs between these two layers is an open design question: durable state that is neither a static instruction nor a single session's transcript. \citet{hu2025memory} and \citet{zhang2026ace}, already cited in \Cref{sec:discuss:future}, motivate an accumulating layer. \citet{packer2024memgpt} reframes the LLM as an operating system with paged memory. \citet{chhikara2025mem0} builds a production-oriented memory store that survives restarts, while \citet{xu2025amem} proposes a research agentic-memory design. \citet{wang2024agentworkflowmemory} captures reusable procedural traces. \citet{shinn2023reflexion} accumulates self-reflection traces via verbal reinforcement across attempts. Surveys by \citet{zhang2024memorysurvey} and \citet{huang2026rethinkingmemory} map candidate mechanisms.

MemGym makes the memory question more concrete. It evaluates dynamic memory formation across tool-use dialogue, deep research, coding, and computer use, rather than only retention of personal facts in chat~\citep{xu2026memgym}. This moves the question from ``where should notes be stored?'' to ``which remembered state helps the agent keep working correctly?''

The same persistence question recurs on the human side. \Cref{sec:discuss:future} already cites longitudinal autonomy evidence from \citet{anthropic2025internal} and \citet{anthropic2026autonomy}. \citet{dellacqua2025cybernetic}'s field experiment with 776 Procter~\& Gamble professionals, together with longitudinal and organisational studies of Copilot rollouts~\citep{stray2025copiloteval} and AI-teamwork trajectories~\citep{xiao2025aiteamwork}, report shifts in human-AI work dynamics as collaboration accumulates. \citet{wang2023voyager} illustrates an embodied agent that accumulates a skill library across tasks. \citet{mollick2024cointelligence} frames the human-AI working relationship as co-intelligence.

Whether a single substrate can carry both a user's personal instruction hierarchy and a shared organisational context while preserving the file-based transparency of CLAUDE.md that \Cref{sec:context} documents is an open architectural question. How session-scoped permissions interact with such a substrate, without reintroducing the resume-restoration concern that \Cref{sec:persist} closes as a deliberate safety choice, is a further open question. The OpenClaw memory subsystem (dreaming, daily notes, hybrid retrieval) and Hermes Agent's WAL-mode SQLite session store with full-text search, both documented in \Cref{sec:compare,sec:compare:hermes}, provide two concrete reference points for how a persistent substrate can sit between a static instruction hierarchy and a single-session transcript.

\subsection{Harness Boundary Evolution: Where, When, What, and with Whom the Agent Acts}
\label{sec:future:boundary}

\Cref{sec:discuss:future} cites \citet{anthropic2026harness}'s argument that better models move the space of useful harness combinations rather than eliminating it. Whether that movement will be most pronounced in \emph{where} the harness runs, \emph{when} it acts, \emph{what} it acts on, or \emph{with whom} it coordinates is not resolved by the source-level analysis in \Cref{sec:arch,sec:turn,sec:auth,sec:ext,sec:context,sec:subagent,sec:persist}. Each of the four has an active research literature that the paper touches only in passing.

\paragraph{Where.} \citet{anthropic2026managed}'s Managed Agents design virtualizes session, harness, and sandbox into independently replaceable interfaces. This extends the virtual-memory analogy that \citet{packer2024memgpt} applies to context-window management and that \citet{karpathy2023llmos} popularizes more broadly. \citet{khattab2024dspy} treats the harness itself as a compile target. 

\paragraph{When.} \Cref{sec:discuss:future} already introduces KAIROS as a feature-gated illustration, motivated by the $+12\%$ to $+18\%$ task-pass gain that \citet{chi2025proactive} report and the sharp preference penalty (47\% vs.\ 80\% to 90\%) restricted to the high-frequency \emph{Persistent Suggest} variant. \citet{liu2025innerthoughts}, \citet{pu2025codellaborator}, and \citet{lee2025sensibleagent} extend the proactivity design space across programming and ambient-interface settings. \citet{pasternak2025probe} and \citet{sun2025userville} introduce benchmarks and training regimes aimed at sharpening it, and \citet{deng2025proactivesurvey} surveys the broader literature.

\paragraph{What.} Vision-language-action work extends the harness beyond textual tool returns. \citet{brohan2023rt2} and \citet{black2024pi0} train VLA policies that execute physical actions, and \citet{ahn2022saycan} grounds plans in robot affordances. Industry systems such as \citet{figure2025helix} and \citet{nvidia2025gr00t} push similar ideas into humanoid control. These systems face the reversibility-weighted risk principle (\Cref{tab:principles}) at a cost asymmetry that the principle names but does not quantify for non-textual actions.

\paragraph{With whom.} Role-differentiated multi-agent systems (\citet{hong2024metagpt}, \citet{li2023camel}, \citet{chen2024agentverse}, \citet{qian2024chatdev}) compose agents with distinct responsibilities. Multi-agent debate~\citep{du2023debate, liang2024divergent} and graph-structured workflows~\citep{zhuge2024gptswarm} explore alternatives to the parent/subagent pattern of \Cref{sec:subagent}. \citet{guo2024massurvey} surveys this space.


Whether a single harness architecture can span all four extensions, or whether the ``harness combinations'' \citet{anthropic2026harness} describes will fragment into specialised stacks, is an open design question. The \emph{when}-extension directly continues the Capability-versus-Adaptability tension in \Cref{tab:tensions}. The \emph{with-whom}-extension partially maps onto Capability-versus-Reliability but raises cross-agent consistency concerns that \Cref{tab:tensions} does not itself cover. The \emph{where}- and \emph{what}-extensions raise further questions the paper's current subsystem boundaries do not cover: which governance obligations attach when harness components become hosted services (\Cref{sec:future:governance}), and how reversibility-weighted risk (\Cref{tab:principles}) scales to physical rather than textual effects. How these extensions compose across axes, rather than within any one, is not something the paper's single-subsystem analyses can resolve.

\subsection{Horizon Scaling: From Session to Scientific Program}
\label{sec:future:horizon}

\Cref{sec:values:five} defines Reliable Execution as spanning ``both single-turn correctness and long-horizon dependability.'' How the architecture documented in \Cref{sec:arch,sec:turn,sec:context,sec:subagent,sec:persist} continues to support long-horizon dependability as autonomous work extends beyond a single session is an open question. Its primary units are the turn, the session, and the sub-agent. A growing literature targets this regime. \citet{lu2024aiscientist} present an end-to-end autonomous research pipeline producing draft manuscripts. \citet{beel2025evalsakana} provide an independent SIGIR Forum evaluation of that pipeline, characterising what ``autonomous research'' currently delivers and where it falls short. \citet{gottweis2025coscientist} develop a multi-agent hypothesis-generation system that runs across days rather than turns, and \citet{novikov2025alphaevolve} pursue algorithmic discovery over timescales that previously took human experts weeks. \citet{kwa2025metrtimehorizon}'s METR study measures the task duration at which frontier agents succeed with fixed reliability, giving an empirical frame for this scaling question.

The Quantitative Goal Persistence benchmark names a narrower failure mode: an agent can make reasonable local tool calls and still stop before a verifier confirms enough valid work units~\citep{cai2026pushagent}. For long-horizon coding agents, progress needs an external notion of done, not only a fluent final message.

Against the paper's analysis, long-horizon deployment tests whether the context-management pipeline of \Cref{sec:context}, the last-assistant-text return policy of \Cref{sec:subagent}, and the append-only persistence of \Cref{sec:persist} remain sufficient when sessions compose into multi-session programs. \Cref{sec:discuss:empirical} already frames this as ``a directly measurable empirical question'' that source-level analysis cannot resolve. Horizon scaling restates that question at the scale of weeks: whether the harness layer alone closes the gap, whether a cross-session memory substrate (\Cref{sec:future:persistence}) is required, or whether horizon-scale work demands coordination primitives beyond session, sub-agent, and memory, is not something the paper's session-scoped analyses can settle.

\tierA{} A development subsequent to the \code{v2.1.88} snapshot analyzed here makes this question concrete. Claude Code \code{v2.1.154} introduced \emph{dynamic workflows}, in which the model writes a JavaScript orchestration script that a background runtime executes, fanning out to many subagents (bounded at up to a thousand per run) while intermediate results live in script variables outside the model's context window rather than passing through it~\citep{claudecode2026workflows}. Released alongside Claude Opus 4.8~\citep{anthropic2026opus48}, this is precisely a coordination primitive beyond the session, sub-agent, and memory units analyzed here: it relocates the orchestration logic itself out of the conversation, extending the \emph{context as scarce resource} and \emph{isolated subagent boundaries} principles (\Cref{tab:principles}) from individual subagents to the orchestration layer. Whether orchestration-as-code becomes the dominant long-horizon primitive, and how its token cost and reduced human mid-run oversight trade against its reliability gains, is an open question the session-scoped analysis here cannot settle.


\subsection{Governance and Oversight at Scale}
\label{sec:future:governance}

Emerging AI regulation adds an external constraint on the architectures that implement the Authority hierarchy of Anthropic, operators, and users documented in \Cref{sec:values:five}. Which logging, transparency, and human-oversight affordances coding-agent architectures should expose under that external constraint remains an open design question. The European Commission's GPAI Code of Practice~\citep{eugpai2025cop} and implementation guidelines~\citep{eugpai2025guidelines} illustrate how general-purpose AI governance is moving toward more explicit expectations for documentation, risk management, transparency, and oversight. The MIT AI Agent Index~\citep{staufer2026agentindex} and the International AI Safety Report~\citep{bengio2026safety}, already cited in \Cref{sec:discuss:future}, motivate the disclosure and oversight side of this constraint. The Bartz v.\ Anthropic ruling~\citep{bartzanthropic2025} adds an input-side constraint on training-data sourcing, distinct from the output-side copyright questions about AI-generated code that emerging cases address separately. An OECD report on AI governance frameworks~\citep{oecd2025governing} and an early analysis of compliance obligations for agent providers by \citet{nannini2026agentlaw} sketch what regulator-facing interfaces might look like without prescribing specifics.

\tierA{} Regulation can also act on a model's availability, not only on what an architecture must disclose. In June 2026 a US export-control order barred foreign nationals from Anthropic's Fable~5 and Mythos~5 models, and within days of launch Anthropic disabled both models for all users worldwide~\citep{anthropic2026fableaccess}. Government policy of this kind can directly determine whether such a model stays deployed at all.

Agentic safety benchmarks also move the safety object from text to environment state. Boiling the Frog evaluates whether multi-turn workspace edits become unsafe under incremental attacks~\citep{bisconti2026boilingfrog}. This matters for coding agents because the risky object may be the modified repository, not the final answer text.

Read against the permission pipeline analyzed in \Cref{sec:auth}, two properties of the current architecture are open under this constraint. First, the deny-first evaluation the paper documents is internally auditable through session transcripts (\Cref{sec:persist}) but not yet externally auditable in the forms that emerging frameworks such as the GPAI Code of Practice~\citep{eugpai2025cop} contemplate. Second, whether the \emph{values-over-rules} principle, which the paper pairs with deterministic guardrails, admits the kind of explicit rule articulation that compliance review may call for is a further open question. Both properties lie within the harness rather than the model, which is where future architectures may need to expose new interfaces.


\subsection{Long-term Developer Capability}
\label{sec:future:lens}

\Cref{sec:values:lens} introduces long-term developer capability as a cross-cutting question rather than a co-equal design value. \Cref{sec:discuss:tensions,sec:discuss:empirical} extend the question with external evidence, including perceived-versus-measured productivity, comprehension loss, complexity accrual, technical-debt persistence, neural-connectivity persistence, and early-career hiring decline. Practitioners say the same: reflecting on his own move to agent-driven coding, \citet{karpathy2026sequoia} warns that ``you can outsource your thinking, but you can't outsource your understanding.'' \Cref{sec:conclude} then pivots: ``Future systems could treat that sustainability gap as a first-class design problem, not a downstream evaluation metric.'' Whether that pivot is possible, and what architectural mechanisms a first-class treatment would require, is the last of the open questions this section records.

This also links to process supervision. Human-LLM collaborative planning work argues that users need process-level control, not only outcome-level review, and organizes steering by mode, scope, and edit level~\citep{he2026steer}. For coding agents, the analogous question is what a user can inspect and change while the work is still unfolding.

Two sub-questions separate the measurement gap from the design gap. First, whether the empirical claims that motivate the question are measurable at session granularity. The existing citations operate at session to multi-month scales (\citet{becker2025measuring}'s 16-developer RCT, \citet{shen2026skill}'s comprehension-test comparison, \citet{kosmyna2025brain}'s EEG study, \citet{he2026cursor}'s 807-repository causal analysis, \citet{liu2026techdebt}'s 304{,}000-commit audit, \citet{rak2025aihiring}'s hiring series), but the harness documented in \Cref{sec:arch,sec:turn,sec:context} exposes no per-session signal for comprehension or convention drift. Related work on programmer interaction modes~\citep{barke2023groundedcopilot} and AI-induced code-security regressions~\citep{perry2023insecurecode} sketches session-granularity measurement, and \citet{aiersilan2026vibecheck} proposes a protocol for session-level cognitive-offloading probes. Second, whether architecture can respond to such measurements once they exist (an analogue of the generator-evaluator separation~\citep{anthropic2026harness} applied to the human loop, comprehension-preserving surfaces, or mechanisms not yet named) is the design-gap question \Cref{sec:conclude} poses. The paper takes no position on which mechanism class is appropriate. It also does not settle whether the harness documented here is the right locus for that action, rather than the IDE, the organisation, or the human development loop. The related work surveyed in \Cref{sec:related} and the sustainability pivot of \Cref{sec:conclude} mark where this paper leaves the question.
